% Volume IX: Institutional Design
% Part XVII. Designing for Delayed Truth
% Chapter 98: Layered Verdicts
% Spine: proof-spines/ch098-spine.md

\chapter{Layered Verdicts}
\label{ch:ch098}

Institutions should not force every inquiry into a binary of true or false when the evidentiary situation is more structured than that. They can instead issue typed determinations such as substantiated, unsubstantiated, disproven, unresolved, procedurally barred, or outside jurisdiction. A layered verdict says what kind of closure has been reached, at what level, and for what purpose.

The crucial requirement is that these statuses remain distinct. ``Not proven'' must not silently become either ``false'' or ``probably true,'' because that collapse lets administrative convenience rewrite epistemic meaning. A mature verdict architecture marks the difference between a claim that has been defeated, a claim that has not yet been established, and a claim the institution is not authorized or prepared to decide.
